\paragraph{We bound sampling and selection error, not construct validity.}
A confidence bound answers ``if I ran this same evaluation on new interventions,
what would I get?''. It does not answer ``is interchange-intervention accuracy
the right thing to measure?''. A hypothesis can score highly for the wrong
reason: subspace patching admits interpretability illusions
\citep{makelov2024subspace}, and faithfulness scores move when the ablation
scheme changes \citep{miller2024faithfulness,zhang2024bestpractices}. Our bounds
inherit whatever the chosen metric means; they make the number reliable, not
necessarily meaningful. The two concerns compose: an illusory hypothesis with a
tight, valid bound is still illusory.

\paragraph{Estimated ablation values.}
Mean ablation replaces a component by an average over a reference distribution.
In our experiments that average is computed once, before any resampling, and held
fixed, so conditional on it the per-instance scores are i.i.d.\ functions of a
single instance and the guarantees apply as stated. A practitioner who recomputes
the ablation values inside each evaluation --- or who uses resample ablation ---
is outside the stated guarantees, which is why the protocol asks for the values to
be frozen.

\paragraph{Few clusters.}
Template-level inference is limited by the number of templates, not the number of
prompts. With the \MyIOINTemplates{} templates we wrote, the cluster-robust band
covers \MyClusterClusterCov{} rather than the nominal $95\%$; our simulations put
the shortfall at about three points for $G\approx 25$ and none for $G\ge 50$
(Table~\ref{tab:cluster}). Few-cluster inference is known to remain
anti-conservative \citep{mnw2023clusterrobust}; the honest reading is that a
template-level guarantee needs more templates, not more prompts.

\paragraph{The bootstrap band is asymptotic.}
Its validity rests on Gaussian approximation for maxima
\citep{chernozhukov2013gaussian,cck2019moment}, which requires variances bounded
away from zero --- a condition that a hypothesis scoring exactly $1$ on every
sampled instance violates. We floor the studentising standard deviation at
$n^{-1/2}$, provide the finite-sample-floored variant of
Section~\ref{sec:methods}(c) for classes where this matters, and verify coverage
empirically in every setting we report. The empirical-Bernstein and Occam bounds
have no such caveat but pay a $\log|\Cclass|$ price.

\paragraph{Continuous hypothesis classes.}
Distributed alignment search optimises a rotation of the residual stream by
gradient descent \citep{geiger2024finding}, so its hypothesis class is
continuous and its search is adaptive. We handle only finite, pre-enumerated
classes directly; for the continuous case the honest options are sample splitting
(which we recommend) or an Occam bound over a discretisation, which requires a
covering-number argument we do not develop here.

\paragraph{Our hypothesis classes are ours, not the field's.}
The magnitude of the optimism we report on GPT-2 depends on the class we
enumerate ($|\Cclass|=\MyIOIm{}$ circuits built from the seven functional head
groups of \citet{wang2023interpretability} plus size-stratified random subsets).
A different class would give different numbers. That dependence is the point ---
optimism is a property of the search, not of the model --- but it means our
specific figures should be read as an illustration of magnitude rather than as a
universal constant.

\paragraph{The template population is a modelling choice.}
Our cluster-robust analysis resamples the \MyIOINTemplates{} templates we wrote.
Those templates are themselves a sample from an ill-defined population of
``sentences that pose an indirect-object-identification problem'', and no
resampling scheme can capture that. We report the cluster-robust interval as a
lower bound on the honest uncertainty about the behaviour, not as its exact
quantification.

\paragraph{Reusable-holdout constants are loose.}
Thresholdout \citep{dwork2015reusable} comes with a generalisation guarantee
whose constants are far too conservative to yield a usable certificate at our
sample sizes. We therefore report it as an empirically validated mechanism, and
recommend a plain held-out set when a guarantee is required.

\paragraph{Scale.}
All experiments run on four tiny transformers and GPT-2 small, within
\MyComputeHours{} GPU-hours on a single consumer accelerator. Nothing in the
statistics depends on model size --- the input is a score matrix --- but we have
not verified that the empirical magnitudes (effective number of hypotheses,
intra-template correlation) transfer to larger models, where circuits are less
sparse and behaviours more heterogeneous.
